This section is independent of the direct-calculus coefficients. It fixes
the entire indexed BFS943 tuple $(\beta,K_0,M_C,M_F)$ of package 133.
The finite recipe in Appendix~\ref{app:constants} and the supplied integer
payloads specify those four numbers completely. In particular
\[
 \beta\ge31\text{ is odd},\quad K_0>0,\quad
 0<M_C<\beta<M_F<2\beta.
\]
The largest numeral has 592,884 bits. The already shifted $M_F$ must not
be shifted a second time or mixed with another coefficient tuple.

\subsection{Complete formulas}

For the 83-operation equation the ten unknowns are
\[
 J,F,\alpha,i,T,w,Y,n,\eta,\zeta\in\Z_{>0}.
\]
For the normalized 80- and 81-operation equations replace the input $Y$
by $s\in\Z_{>0}$ and define $Y=q^3s$. The common polynomial definitions are
\begin{align*}
 q&=\beta J+1,& h&=q-1,& C&=q-F-\alpha,\\
 X&=wq,&\mu&=q^2+1,& k&=\eta+\zeta,\\
 c&=kY+\eta,& A&=\mu Y(X+1)+1,\\
 \Delta&=A^2-1,& H&=4A-5,\\
 v&=X\mu^2Y^2,&E&=1+v(v+1)k^2,\\
 R&=[h(q-F)+\alpha](q^2-1)+(M_C+qM_F)J,\\
 e&=k-R-1,&B&=(K_0+w)C-F,\\
 U&=ic^2,&a&=\Delta U^2E,&F_0&=1+a,\\
 G&=EF_0,\\
 F_*&=1+4aF_0,&Q&=\Delta(F_*-1),\\
 b&=cT-R,&V&=F_*b-c,&K_+&=X(3c-2X)+2.
\end{align*}
The abbreviation $F_*$ denotes a polynomial, not an additional unknown.
Take the transport weight $L=V^2$ for the 83- and 80-operation forms,
and $L=bV$ for the 81-operation form. For the unnormalized form set
\begin{align*}
 P&=q^3h,&M&=P\mu^2H,\\
 Z&=P\mu^2[H(n-LB)+K_+]
       +H[\mu^2(hY-q^3B)-Pe].
\end{align*}
For the normalized forms set
\[
 M=h\mu^2H,\qquad
 Z=h\mu^2[H(n-LB)+K_+]-H[\mu^2B+he].
\]
In each case the complete equation is
\begin{equation}\label{eq:compact}
 \boxed{Q\bigl(Z^2-M^2GV^2\bigr)+(M^2G-Z^2)=0.}
\end{equation}
It is identically $(Q-1)Z^2-M^2G(QV^2-1)=0$.
The strict positivity of the ten inputs is part of the domain.

\subsection{Root extraction and the free auxiliary coordinate}

The package-133 kernel uses Pell doubling:
\[
 F_*=(1+2a)^2,\qquad Q=4\Delta^2U^2G,\qquad\gcd(E,F_0)=1.
\]
Positive inputs give $q\ge32$, $c\ge3$, $a>c$, $Q>c^2$, and $M>0$.
If $b=0$, then $V=-c$; otherwise $|V|\ge F_*-c>1$. At a zero, put
$z=Z/M$. Since $G\mid Q$, both $Q-1$ and $QV^2-1$ are $-1$ modulo $G$.
Prime valuations in
\[
 (Q-1)z^2=G(QV^2-1)
\]
make $G=g^2$ for an integer $g>0$. Coprimality gives
$E=\tau^2$, $F_0=f_0^2$, and $g=\tau f_0$.

The positive rational $y=|z|/g$ satisfies
\begin{equation}\label{eq:adjacent}
 QV^2-(Q-1)y^2=1.
\end{equation}
It is integral. Indeed, if $y=r/m$ is reduced, then $m^2\mid Q-1$.
For $v_1=|V|>1$, the norm-preserving transformation
\[
 v_2=(2Q-1)v_1-2(Q-1)y,\quad y_2=(2Q-1)y-2Qv_1
\]
keeps $v_2$ integral, preserves the reduced denominator of $y_2$, and
has $-v_1<v_2<v_1$. The preserved norm excludes $v_2=0$. Descending in $|v_1|$ reaches one, where $y=1$.
Consequently $z=\epsilon gy$ for $\epsilon=\pm1$.

Define, in the proof only, $f=1+2a$ and $I=2Ug$. Then
\[
 f^2-\Delta I^2=1,\quad f^2=F_*,\quad Q=(\Delta I)^2,
 \quad c^2\mid I.
\]
The retained main-folding theorem permits either sign of $V$ and yields
$c=Y_p(A)$ for odd $p\ge1$. The first norm gives
$k=2Y_{n_0}(2v+1)$. The case $n_0=1$ would give $c=2Y+1$ strictly
between the first two Pell ordinates, so $n_0\ge2$. In particular
\[
 k\ge8v+4,\qquad c>k>8q^5,\qquad K_+>0.
\]

The new improvement is that $b=cT-R$ needs no preliminary shift
$T+\alpha$ when the transport argument is proved with absolute values.
At a zero, $b$ cannot vanish: otherwise
\[
 y^2=c^2+\frac{c^2-1}{Q-1}
\]
lies strictly between consecutive integers. For nonzero $b$, its sign
equals that of $V$ and $|b|<|V|$. Also
\[
 a^2-g^2=E[\Delta U^2E(\Delta U^2-1)-1]>0,
\]
so $g<a$ and $|V|\ge F_*-c>3g$. From \eqref{eq:adjacent},
$y<\tfrac32|V|$, hence $gy<V^2/2$.

Put $\kappa=K_+/H>0$, $t=B/h$, $\rho=e/\mu^2$, and
$W_*=1+hL$. The exact carrier identities are
\[
 \frac ZM=n+\kappa+Y/q^3-W_*t-\rho
 \quad\hbox{or}\quad
 \frac ZM=n+\kappa-W_*t-\rho.
\]
The one-sided estimate
\[
 b-e=cT-k+1>0,\qquad
 \rho<\frac b{\mu^2}\le\frac{|b|}{\mu^2}<\frac{|V|}{\mu^2}
\]
uses only $T\ge1$ and $c>k$. If $B\le-1$, then for $L=V^2$ either
carrier gives
\[
 n<\frac{V^2}{2}-V^2+\frac{|V|}{\mu^2}<0,
\]
a contradiction. Thus $B\ge0$, and the definition of $B$ forces $C>0$.
This proves the essential new soundness step for the 83- and 80-operation
equations. It never assumes the deleted shift has a positive inverse.

\subsection{Why the smaller weight works after normalization}

For $L=bV$ a stronger bound is needed. The literal identities $q\mid X,Y$
in the normalized chart exclude the exceptional main index $p=3$ before
transport. Suppose $p=3$, so $c=4A^2-1$. A first index $n_0\ge3$ would
give $kY>32v^2Y>4A^2>c$, so $n_0=2$ and $k=8v+4$.
Writing $D=(X+1)^2-2XY$ then gives the exact identity
\[
 \eta=4\mu^2Y^2D+8\mu Y(X+1)+3-4Y.
\]
The inequalities $0<\eta<k$ exclude $D\ge2X$ and $D\le0$.
For the latter, $D=1+X(X+2-2Y)\le-1$ implies $X+1<2Y$, making the
display negative. Thus $0<D<2X$, which forces
$2Y-X\in\{1,2\}$. This contradicts $q\mid X,Y$ and $q\ge32$.
Hence $p\ge5$ and $c>4\Delta+3$.

In \eqref{eq:adjacent}, put $S=\Delta I$ and $Q=S^2$. Pell classification
gives an odd auxiliary index; its first and third absolute $V$ values
are $1$ and $4Q-3$. The first is already excluded. The second is
impossible because $V=F_*b-c$ would give either
\[
 F_*(b-4\Delta)=c-4\Delta-3,
 \qquad F_*(b+4\Delta)=c+4\Delta+3,
\]
whose right side is strictly between zero and $F_*$. Therefore
$|V|\ge16Q^2-20Q+5>8Q^2$. As $c<F_*$,
\[
 |b|>8Q^2/F_*-1>4\Delta^3I^2-1>4g.
\]
Now $L=bV=|b||V|>0$, $gy<3L/8<L/2$, and
$\rho<|b|/\mu^2<L/2$. If $B\le-1$, the carrier again gives $n<0$.
This proves the missing sign step for the 81-operation equation.

\subsection{Later decoding, converse, and exact size}

After $C>0$, the retained mask bound gives $3q<R<q^4-2$. The signed-index
lemma uses the actual residues $V\equiv-c\pmod f$ and
$V\equiv-R\pmod c$, with $c^2\mid I$; it allows either sign of $V$ and
contains no condition $T>\alpha$. It gives $p=R$. The original parity,
denominator, first-index, and whole-product projector arguments then
restore the native source and its finite refutation. The normalized
carrier supplies a stronger denominator premise; no nonunit is cancelled
modulo a composite modulus.

Conversely the parent's canonical witnesses give a positive old
$T_{\rm old}$ and $q^3\mid Y$. Set $T=T_{\rm old}+\alpha$. For the
normalized equations also set $s=Y/q^3$ and
$n=n_{\rm old}+s$. These are positive. The exact substitutions give
\[
 Z_{\rm old}=q^3Z_{\rm new},\quad
 M_{\rm old}=q^3M_{\rm new},\quad
 N_{\rm old}=q^6N_{\rm new},
\]
so the new polynomial vanishes. Soundness was proved directly because
the inverse coordinate shift need not preserve positivity.

The three live operation counts are respectively
$(50,18,15)$, $(49,18,14)$, and $(48,18,14)$ for multiplication,
addition, and subtraction. Put $q_0=\beta J$ and $k_0=\eta+\zeta$.
Their highest homogeneous forms are
\begin{align*}
 P_{83}^{\rm top}&=2^{14}q_0^{200}w^{46}Y^{104}i^{20}k_0^{60}
                  (q_0-F)^4(q_0-F-\alpha)^2,\\
 P_{81}^{\rm top}&=-2^{10}q_0^{384}w^{34}s^{80}i^{14}k_0^{46}T^2,\\
 P_{80}^{\rm top}&=2^{14}q_0^{506}w^{46}s^{108}i^{20}k_0^{64}T^4
                  (q_0-F-\alpha)^2.
\end{align*}
They are nonzero for the authentic tuple and have degrees 436, 560,
and 750. Global gate-by-gate upper bounds and full univariate polynomial
replays independently verify attainment. A finite ray calculation alone
would not supply the global upper bound; both are included.
